\ifdefined\gkver\else\def\gkver{student}\fi
\documentclass[\gkver]{gaokaozhenti}
\begin{document}

\chapter{2026年海南}

\begin{enumerate}
\item
%% number: 1
%% paperName: 2026年海南高考第 1 题 3 分
%% typeId: 1
%% type: 单选题
%% chapter: 功和能
%% point: 动能势能
%% method: 无
%% score: 3
%% degree: 950
%% duplicateId: 0
%% body:
在 2026 年世界乒乓球团体锦标赛中，我国乒乓球队男队、女队双双卫冕成功。比赛中，乒乓球从水平台位置 1 经最高点位置 2 落到球台位置 3，运动轨迹如图所示，则\xzanswer{C}
\onepicture[w=6.5cm]{figs/01.png}
\fourchoices[answer=C]
{从位置 1 到位置 2 乒乓球的重力做正功}
{从位置 2 到位置 3 乒乓球的重力做负功}
{从位置 1 到位置 2 乒乓球的重力势能增大}
{从位置 1 到位置 3 乒乓球的重力势能减小}

%% answer: C
%% memo:
\memoanswer{
A、C．从位置 1 到位置 2 乒乓球的高度增加，则重力做负功，乒乓球的重力势能增大，A 错误，C 正确；

B．从位置 2 到位置 3 乒乓球的高度降低，则重力做正功，B 错误；

D．从位置 1 到位置 3 乒乓球的高度先增加后减小，则重力势能先增加后减小，D 错误。

故选 C。
}

\item
%% number: 2
%% paperName: 2026年海南高考第 2 题 3 分
%% typeId: 1
%% type: 单选题
%% chapter: 原子和原子核
%% point: 天然放射性
%% method: 无
%% score: 3
%% degree: 950
%% duplicateId: 0
%% body:
放射性同位素碘 131 常作为放射性示踪剂，用于甲状腺病的诊断。碘 131 发生 $\beta$ 衰变的方程为 \ce{^{131}_{53}I -> ^{m}_{n}X + ^{0}_{-1}e}，则\xzanswer{A}
\fourchoices[answer=A]
{$m = 131$，$n = 54$}
{$m = 130$，$n = 53$}
{$m = 131$，$n = 52$}
{$m = 132$，$n = 53$}

%% answer: A
%% memo:
\memoanswer{
由质量数守恒得 $131 = m + 0$，解得 $m = 131$；

由电荷数守恒得 $53 = n + (-1)$，解得 $n = 54$。

故选 A。
}

\item
%% number: 3
%% paperName: 2026年海南高考第 3 题 3 分
%% typeId: 1
%% type: 单选题
%% chapter: 光学
%% point: 光电效应
%% method: 无
%% score: 3
%% degree: 850
%% duplicateId: 0
%% body:
现有 a、b 两束单色光，a 光的频率小于 b 光的频率。用 a 光照某金属表面能发生光电效应。下列说法正确的是\xzanswer{A}
\fourchoices[answer=A]
{若用 b 光照射该金属，则能发生光电效应}
{若用 b 光照射该金属，则该金属的逸出功减小}
{若增大 a 光强度，则该金属的逸出功增大}
{若减弱 a 光强度，则不能发生光电效应}

%% answer: A
%% memo:
\memoanswer{
A．已知 $\nu_{a} < \nu_{b}$，且 a 光能发生光电效应，说明 $\nu_{a} > \nu_{0}$，因此 $\nu_{b} > \nu_{a} > \nu_{0}$，b 光频率也大于金属极限频率，能发生光电效应，故 A 正确；

B．逸出功是金属的固有属性，与入射光频率无关，用 b 光照射不会改变金属逸出功，故 B 错误；

C．逸出功与入射光强度无关，增大 a 光强度不会改变金属逸出功，故 C 错误；

D．光电效应的发生仅由入射光频率决定，与光强无关，只要 a 光频率大于极限频率，即使减弱光强，依然能发生光电效应，故 D 错误。

故选 A。
}

\item
%% number: 4
%% paperName: 2026年海南高考第 4 题 3 分
%% typeId: 1
%% type: 单选题
%% chapter: 运动和力
%% point: 动量守恒定律
%% method: 无
%% score: 3
%% degree: 650
%% duplicateId: 0
%% body:
如图，在光滑水平面上，a 球以初速度 $v$ 水平向右运动，与前方静止的 b 球发生弹性正碰。若 a 球的质量 $M$ 远大于 b 球的质量 $m$，则碰撞后\xzanswer{B}
\onepicture[w=6.5cm]{\includesvg{figs/04}}
\fourchoices[answer=B]
{a 球的速度为 0}
{a 球的速度略小于 $v$}
{两球以相同的速度向右运动}
{两球的总动量比碰撞前的小}

%% answer: B
%% memo:
\memoanswer{
设向右为正，因两球发生弹性正碰，则碰后两球的总动量等于碰前的总动量，由动量守恒和能量守恒关系可知 $Mv = Mv_{1} + mv_{2}$，$\frac{1}{2}Mv^{2} = \frac{1}{2}Mv_{1}^{2} + \frac{1}{2}mv_{2}^{2}$，

解得 $v_{1} = \frac{M - m}{M + m}v$，$v_{2} = \frac{2M}{M + m}v$，

因 $M \gg m$ 可知 $v \approx v_{1} > 0$，$v_{2} > v_{1}$，

即 a 球的速度略小于 $v$，b 球的速度大于 a 球的速度。

故选 B。
}

\item
%% number: 5
%% paperName: 2026年海南高考第 5 题 3 分
%% typeId: 1
%% type: 单选题
%% chapter: 光学
%% point: 光的折射
%% method: 无
%% score: 3
%% degree: 650
%% duplicateId: 0
%% body:
如图，一束单色光由空气射向某一透明介质，入射角为 $60^{\circ}$，反射光线与折射光线的夹角为 $75^{\circ}$，则该介质的折射率为\xzanswer{B}
\onepicture[w=4.8cm]{figs/05.png}
\fourchoices[answer=B]
{$\sqrt{2}$}
{$\frac{\sqrt{6}}{2}$}
{$\sqrt{3}$}
{$\frac{4}{3}$}

%% answer: B
%% memo:
\memoanswer{
已知入射角为 $60^{\circ}$，因此反射角也为 $60^{\circ}$，即反射光线与法线的夹角为 $60^{\circ}$；

由几何关系可知：反射光线与界面的夹角为 $90^{\circ}-60^{\circ}=30^{\circ}$；

因此折射光线与界面的夹角为 $75^{\circ}-30^{\circ}=45^{\circ}$；

由此可得折射角（折射光线与法线的夹角）$r=90^{\circ}-45^{\circ}=45^{\circ}$；

根据光的折射定律 $n=\frac{\sin 60^{\circ}}{\sin 45^{\circ}}=\frac{\sqrt{6}}{2}$。

故选 B。
}

\item
%% number: 6
%% paperName: 2026年海南高考第 6 题 3 分
%% typeId: 1
%% type: 单选题
%% chapter: 电磁感应
%% point: 远距离输电
%% method: 无
%% score: 3
%% degree: 650
%% duplicateId: 0
%% body:
2025 年 12 月，海南自贸港重大项目——海南 500 千伏主网架工程全面投产。海南电网主网架首次实现电压等级从 $220 \UkV$ 到 $500 \UkV$ 的历史性跨越。若保持输送电功率和输电线电阻都不变，采用 $500 \UkV$ 输电与采用 $220 \UkV$ 输电在输电线上损失的电功率之比为\xzanswer{D}
\fourchoices[answer=D]
{25:11}
{11:25}
{625:121}
{121:625}

%% answer: D
%% memo:
\memoanswer{
输电过程中，输送功率 $P$、输电电压 $U$、输电电流 $I$ 满足 $P = UI$，因此输电电流 $I = \frac{P}{U}$；

输电线上损失的功率为电流的热功率，即 $\Delta P = I^{2}R$，联立得 $\Delta P = \frac{P^{2}R}{U^{2}}$；

两种输电方式损失功率之比为 $\frac{\Delta P_{500}}{\Delta P_{220}} = \frac{U_{220}^{2}}{U_{500}^{2}} = \frac{220^{2}}{500^{2}} = \frac{121}{625}$。

故选 D。
}

\item
%% number: 7
%% paperName: 2026年海南高考第 7 题 3 分
%% typeId: 1
%% type: 单选题
%% chapter: 曲线运动
%% point: 人造地球卫星
%% method: 无
%% score: 3
%% degree: 450
%% duplicateId: 0
%% body:
2026 年 4 月，长征八号运载火箭在海南文昌成功发射，以“一箭十八星”方式将卫星送入预定轨道。若某卫星绕地球做圆周运动的轨道半径为 $r$，地球半径为 $R$，地球表面重力加速度为 $g$，忽略地球自转影响，则该卫星的\xzanswer{B}
\fourchoices[answer=B]
{线速度大小为 $\sqrt{gR}$}
{加速度大小为 $\frac{gR^{2}}{r^{2}}$}
{周期为 $2\pi\sqrt{\frac{R^{2}}{gr^{2}}}$}
{角速度大小为 $\sqrt{\frac{g}{R}}$}

%% answer: B
%% memo:
\memoanswer{
A．根据题意，对于地球表面的物体，由万有引力等于重力有 $\frac{GMm}{R^{2}} = mg$，可得 $GM = gR^{2}$；

由万有引力提供向心力有 $\frac{GMm}{r^{2}} = m\frac{v^{2}}{r}$，联立解得 $v = \sqrt{\frac{gR^{2}}{r}}$，故 A 错误；

B．由万有引力提供向心力有 $\frac{GMm}{r^{2}} = ma$，解得 $a = \frac{GM}{r^{2}}$，代入 $GM = gR^{2}$，可得 $a = \frac{gR^{2}}{r^{2}}$，故 B 正确；

C．由万有引力提供向心力有 $\frac{GMm}{r^{2}} = m\frac{4\pi^{2}}{T^{2}}r$，解得 $T = \sqrt{\frac{4\pi^{2}r^{3}}{GM}}$，代入 $GM = gR^{2}$，可得 $T = 2\pi\sqrt{\frac{r^{3}}{gR^{2}}}$，故 C 错误；

D．由万有引力提供向心力有 $\frac{GMm}{r^{2}} = m\omega^{2}r$，解得 $\omega = \sqrt{\frac{GM}{r^{3}}}$，代入 $GM = gR^{2}$，可得 $\omega = \sqrt{\frac{gR^{2}}{r^{3}}}$，故 D 错误。

故选 B。
}

\item
%% number: 8
%% paperName: 2026年海南高考第 8 题 3 分
%% typeId: 1
%% type: 单选题
%% chapter: 力物体的平衡
%% point: 动态平衡
%% method: 无
%% score: 3
%% degree: 450
%% duplicateId: 0
%% body:
如图，两个木块 A 和 B 叠放在水平地面上，一根轻质弹性绳上端固定于天花板上的 $O$ 点，下端跨过一个固定在 $P$ 点的光滑钉子斜拉着木块 A，$P$ 点在 $O$ 点的正下方。已知弹性绳原长小于线段 $OP$ 的长度，A 和 B 均处于静止状态，弹性绳的弹力大小与伸长量的关系满足胡克定律，A 和 B 均可视为质点，下列说法正确的是\xzanswer{D}
\onepicture[w=6.2cm]{figs/08.png}
\fourchoices[answer=D]
{地面对 A 的摩擦力方向水平向左}
{若 B 的质量增大，则 A 与地面之间的摩擦力变大}
{若将 A 和 B 整体向右平移一小段距离后重新平衡，则 A 对地面的压力变小}
{若将 A 和 B 整体向右平移一小段距离后重新平衡，则 A 对地面的压力变大}

%% answer: D
%% memo:
\memoanswer{
A．绳对 A 的拉力方向是斜向左上方，拉力的水平分量向左，A 有向左运动的趋势，因此地面对 A 的静摩擦力方向水平向右，故 A 错误；

B．对 A、B 整体受力分析，水平方向静摩擦力大小等于弹性绳拉力的水平分量，和 B 的质量没有关系，因此 B 质量增大时，A 与地面的摩擦力不变，故 B 错误；

C、D．设 $PA$ 段绳长为 $L$、竖直高度为 $h$，弹性绳劲度系数为 $k$，对 A、B 整体竖直方向受力平衡，设绳拉力为 $F$，拉力与水平方向夹角为 $\theta$，则 $\sin\theta=\frac{h}{L}$；

总伸长量对应的弹力 $F = k(OP + L - l_{0})$；

竖直方向 $N + F \sin \theta = (m_{\mathrm{A}} + m_{\mathrm{B}}) g$；

整理得 $N = (m_{\mathrm{A}} + m_{\mathrm{B}}) g - k (OP + L - l_{0}) \cdot \frac{h}{L} = (m_{\mathrm{A}} + m_{\mathrm{B}}) g - kh \cdot \frac{OP - l_{0}}{L} - kh$；

当整体向右平移小段距离，$L$ 增大，$N$ 变大，即地面对 A 的支持力变大，根据牛顿第三定律，A 对地面的压力变大，故 D 正确，C 错误。

故选 D。
}

\item
%% number: 9
%% paperName: 2026年海南高考第 9 题 0 分
%% typeId: 2
%% type: 多选题
%% chapter: 机械振动
%% point: 振动图象
%% method: 无
%% score: 0
%% degree: 650
%% duplicateId: 0
%% body:
某质点做简谐运动的振动图像如图所示，则\xzanswer{AC}
\onepicture[w=5.5cm]{\includesvg{figs/09}}
\fourchoices[answer=AC]
{质点的振动周期为 $8 \Us$}
{$t = 2 \Us$ 时质点所受的回复力为 0}
{$t = 3 \Us$ 时质点向平衡位置运动}
{$t = 4 \Us$ 时质点速度为 0}

%% answer: AC
%% memo:
\memoanswer{
A．根据题图可知质点的振动周期为 $8 \Us$，故 A 正确；

B．根据题图可知 $t = 2 \Us$ 时质点处于正向位移最大处，则所受的回复力最大，故 B 错误；

C．根据题图可知 $t = 3 \Us$ 时质点向平衡位置运动，故 C 正确；

D．根据题图可知 $t = 4 \Us$ 时质点处于平衡位置，速度最大，故 D 错误。

故选 AC。
}

\item
%% number: 10
%% paperName: 2026年海南高考第 10 题 4 分
%% typeId: 2
%% type: 多选题
%% chapter: 气体性质
%% point: 气体图象
%% method: 无
%% score: 4
%% degree: 550
%% duplicateId: 0
%% body:
下列图像中，$p$ 表示气体的压强，$V$ 表示气体的体积。对一定质量的理想气体，其中状态 a 的温度大于状态 b 的是\xzanswer{BCD}
\fourchoices[answer=BCD,ispicture=true,h=2.6cm]
{figs/10a.png}
{figs/10b.png}
{figs/10c.png}
{figs/10d.png}

%% answer: BCD
%% memo:
\memoanswer{
A．根据题图可知状态 a 的压强等于状态 b 的压强，状态 a 的体积小于状态 b 的体积，根据 $\frac{pV}{T}=C$ 可知状态 a 的温度小于状态 b 的温度，故 A 错误；

B．根据题图可知状态 a 的压强大于状态 b 的压强，状态 a 的体积等于状态 b 的体积，根据 $\frac{pV}{T}=C$ 可知状态 a 的温度大于状态 b 的温度，故 B 正确；

C．根据 $\frac{pV}{T}=C$ 可得 $p=CT\cdot\frac{1}{V}$，根据题图可知状态 a 与原点的连线斜率大于状态 b 与原点的连线斜率，则状态 a 的温度大于状态 b 的温度，故 C 正确；

D．根据 $\frac{pV}{T}=C$ 可得 $p=CT\cdot\frac{1}{V}$，根据题图可知状态 a 与原点的连线斜率大于状态 b 与原点的连线斜率，则状态 a 的温度大于状态 b 的温度，故 D 正确。

故选 BCD。
}

\item
%% number: 11
%% paperName: 2026年海南高考第 11 题 4 分
%% typeId: 2
%% type: 多选题
%% chapter: 光学
%% point: 光的干涉
%% method: 无
%% score: 4
%% degree: 650
%% duplicateId: 0
%% body:
如图，在双缝干涉实验中，用某一单色光照射双缝，在光屏上观察到干涉条纹，下列说法正确的是\xzanswer{BD}
\onepicture[w=4.5cm]{figs/11.png}
\fourchoices[answer=BD]
{若增大光屏与双缝间的距离，则相邻亮条纹中心的间距变小}
{若增大光屏与双缝间的距离，则相邻亮条纹中心的间距变大}
{若用波长更长的单色光照射，则相邻亮条纹中心的间距变小}
{若用波长更长的单色光照射，则相邻亮条纹中心的间距变大}

%% answer: BD
%% memo:
\memoanswer{
根据 $\Delta x=\frac{L}{d}\lambda$ 分析（其中 $L$ 是光屏与双缝间的距离，$d$ 是双缝间的距离）。

A、B．根据 $\Delta x=\frac{L}{d}\lambda$ 可知若增大光屏与双缝间的距离，则相邻亮条纹中心的间距变大，故 A 错误，B 正确；

C、D．根据 $\Delta x=\frac{L}{d}\lambda$ 可知若用波长更长的单色光照射，则相邻亮条纹中心的间距变大，故 C 错误，D 正确。

故选 BD。
}

\item
%% number: 12
%% paperName: 2026年海南高考第 12 题 4 分
%% typeId: 2
%% type: 多选题
%% chapter: 电场
%% point: 等势面
%% method: 无
%% score: 4
%% degree: 450
%% duplicateId: 0
%% body:
如图，$a$、$b$、$c$、$d$ 是圆心为 $O$ 的圆周的四等分点，$M$、$N$ 分别为 $ab$、$bc$ 的中点。在 $a$、$c$ 分别固定一正点电荷 $+q$，在 $b$、$d$ 分别固定一负点电荷 $-q$，则\xzanswer{AD}
\onepicture[w=5.5cm]{\includesvg{figs/12}}
\fourchoices[answer=AD]
{$O$ 点的电场强度为 0}
{$O$ 点与 $M$ 点的电场强度大小相等}
{$O$ 点的电势比 $M$ 点的低}
{$M$、$N$ 两点的电势相等}

%% answer: AD
%% memo:
\memoanswer{
A．根据题意，由点电荷场强公式 $E = \frac{kQ}{r^{2}}$，由对称性可知，a、c 固定的点电荷在 O 点产生的电场等大反向，b、d 固定的点电荷在 O 点产生的电场等大反向，则 O 点的电场强度为 0，故 A 正确；

B．根据题意，由点电荷场强公式 $E = \frac{kQ}{r^{2}}$ 可知，各点电荷在 M 点产生的电场，如图所示。

\onepicture[w=5.0cm]{figs/12_an1.png}

由几何关系可得，M 点的电场强度大小为 $E_{M} = E_{a} + E_{b} - 2E_{c}\cos\alpha > 0$，则 O 点与 M 点的电场强度大小不相等，故 B 错误；

C、D．根据题意，如图所示。

\onepicture[w=5.0cm]{figs/12_an2.png}

由几何关系可知，过 $OM$ 的直线是 a、b 固定的等量异种点电荷的中垂线，也是 c、d 固定的等量异种点电荷的中垂线，则 $OM$ 是一条等势线；同理，过 $ON$ 的直线是 c、b 固定的等量异种点电荷的中垂线，也是 a、d 固定的等量异种点电荷的中垂线，则 $ON$ 是一条等势线，则 O、M、N 三点的电势相等，故 C 错误，D 正确。

故选 AD。
}

\item
%% number: 13
%% paperName: 2026年海南高考第 13 题 4 分
%% typeId: 2
%% type: 多选题
%% chapter: 电磁感应
%% point: 电磁感应中的变加速
%% method: 无
%% score: 4
%% degree: 350
%% duplicateId: 0
%% body:
如图，两根平行光滑金属导轨固定在同一水平面内，左端接一定值电阻 $R$。$Ox$ 轴平行于金属导轨，在 $1 \Um \leq x \leq 2 \Um$ 的空间区域内存在垂直纸面向里的磁场，磁感应强度的大小 $B$ 随坐标 $x$ 的分布规律为 $B = 0.2x(\UT)$。金属细棒 $ab$ 在水平向右的外力 $F$ 作用下平行于导轨运动，运动过程中始终与导轨垂直并接触良好。已知导轨间距为 $1 \Um$，电阻 $R$ 的阻值为 $0.5 \UO$，金属棒从 $x_{1} = 1 \Um$ 运动到 $x_{2} = 2 \Um$ 的过程中，电阻 $R$ 的电功率保持不变，产生的焦耳热为 $0.6 \UJ$。不计导轨和金属棒的电阻，则金属棒\xzanswer{BC}
\onepicture[w=9.0cm]{\includesvg{figs/13}}
\fourchoices[answer=BC]
{从 $x_{1}$ 到 $x_{2}$ 的过程做匀减速直线运动}
{从 $x_{1}$ 到 $x_{2}$ 的过程受到的外力 $F$ 逐渐增大}
{经过 $x_{1}$ 处时的速度大小为 $5 \Ums$}
{从 $x_{1}$ 到 $x_{2}$ 的过程所用的时间为 $0.2 \Us$}

%% answer: BC
%% memo:
\memoanswer{
A．动生电动势 $E = BLv = 0.2xv \UV$，电阻 R 的电功率保持不变，由 $P=\frac{E^{2}}{R}$ 可知感应电动势 $E$ 恒定，则速度 $v$ 与位移 $x$ 成反比，不是匀减速直线运动，故 A 错误；

C、D．电流恒定则 $xv = k$ 为定值，取过程中极小段时间 $\Delta t$，则有 $\Delta t = \frac{\Delta x}{v} = \frac{x\Delta x}{k}$，则有 $\sum\Delta t=\frac{1}{k}\sum x\Delta x$，由数学知识知 $\sum x\Delta x=(x_{1}+x_{2})(x_{2}-x_{1})\times\frac{1}{2}$，整理可得 $t=\frac{3}{2xv} \Us$；

已知总焦耳热 $Q=I^{2}Rt=\frac{0.2^{2}x^{2}v^{2}}{R}t=0.6 \UJ$，整理可得 $xv = 5 \Umqs$，$t = 0.3 \Us$；

代入 $x_{1} = 1 \Um$，可得 $v_{1} = 5 \Ums$，故 C 正确，D 错误；

B．由于 $xv = k$ 为定值，由数学知识 $(xv)' = x'v + v'x = 0 \Rightarrow v^{2} + ax = 0$，由于位移增大而速度减小，加速度大小减小；

安培力 $F_{\text{安}} = BIL = \frac{BLE}{R} = 0.4x \UN$，导轨光滑，由牛顿第二定律 $F - F_{\text{安}} = ma$，由于加速度为负且大小减小，可推出受到的外力 $F$ 逐渐增大，故 B 正确。

故选 BC。
}

\item
%% number: 14
%% paperName: 2026年海南高考第 14 题 8 分
%% typeId: 4
%% type: 实验题
%% chapter: 气体性质
%% point: 用单分子油膜估测分子大小
%% method: 无
%% score: 8
%% degree: 650
%% duplicateId: 0
%% body:
某同学用油膜法估测油酸分子的大小，实验器材有：油酸、酒精、注射器、烧杯、量筒、浅盘、痱子粉、带有坐标方格的玻璃板、记号笔等。实验步骤如下：

Ⅰ．配制油酸酒精溶液；

Ⅱ．用注射器吸取配制好的油酸酒精溶液，测量 $1 \UmL$ 油酸酒精溶液的总滴数；

Ⅲ．在浅盘里盛上水，将痱子粉均匀地撒在水面上；

Ⅳ．用注射器向水面上滴 1 滴油酸酒精溶液，待油膜形状稳定后，将玻璃板盖在浅盘上，在玻璃板上描下油膜轮廓。

回答下列问题：
\begin{enumerate}
	\item
	本实验将油酸分子简化为球形，并认为形成油膜的油酸分子是 \tkanswer{单}（选填“单”或“多”）层紧密排布的。
	\item
	实验中痱子粉的作用是 \tkanswer{B}（单选，填正确选项的标号）
\threechoices[answer=B]
{促进油酸溶解}
{使油膜轮廓清晰可见}
{增大油膜的面积}
	\item
	实验中，每 $500 \UmL$ 油酸酒精溶液中有纯油酸 $1 \UmL$，每 $1 \UmL$ 这样的溶液有 200 滴，把 1 滴这样的溶液滴入盛水的浅盘中，油膜形状稳定后轮廓如图所示，图中一个小方格的面积为 $1 \Ucmq$。该同学数得图中轮廓内完整的小方格有 106 个，超过半格但不足整格的有 24 个，则该同学测得的油膜面积为 \tkanswer{130} $\Ucmq$。根据实验数据，该同学测得油酸分子的直径为 \tkanswer{$8\times10^{-10}$} $\Um$（保留 1 位有效数字）。
	\item
	若油酸未完全散开就描绘油膜轮廓，导致测得的油膜面积比实际小，则测得的油酸分子直径 \tkanswer{偏大}（选填“偏大”或“偏小”）。
\end{enumerate}
\onepicture[w=6.0cm, align=right]{figs/14.png}

%% answer:
\jdanswer{
\begin{enumerate}
	\item
	单

	\item
	B

	\item
	130，$8\times10^{-10}$

	\item
	偏大
\end{enumerate}
}
%% memo:
\memoanswer{
（1）油膜法估测分子大小的核心模型是假设油酸在水面上充分展开后，形成单层的油酸分子油膜，通过油膜厚度等效为油酸分子的直径。

（2）痱子粉本身不与油酸反应，会漂浮在水面上，油酸扩散时会将痱子粉向四周推开，从而清晰显现出油膜的边界轮廓。油酸本身不溶于水，痱子粉无法促进溶解、也不能改变油膜的实际面积。

故选 B。

（3）油膜面积的计数规则是：完整方格全算，超过半格的按 1 格计、不足半格的舍去，因此总方格数为 $106+24=130$，对应面积 $S = 130 \times 1 \Ucmq = 130 \Ucmq$。

1 滴溶液中纯油酸的体积：$V=\frac{1 \UmL}{500}\times\frac{1}{200}=1\times10^{-5} \UmL=1\times10^{-11} \Umc$。

分子直径 $d=\frac{V}{S}=\frac{1\times10^{-11}}{1.3\times10^{-2}} \Um \approx 8\times10^{-10} \Um$。

（4）分子直径计算公式为 $d=\frac{V}{S}$，若油膜未完全散开，测得的面积 $S$ 偏小，计算得到的直径 $d$ 会偏大。
}

\item
%% number: 15
%% paperName: 2026年海南高考第 15 题 19 分
%% typeId: 4
%% type: 实验题
%% chapter: 电路
%% point: 测量电阻率
%% method: 极值
%% score: 19
%% degree: 450
%% duplicateId: 0
%% body:
某实验小组测量一合金金属丝的电阻率，实验器材有：电源（电动势 $E = 3 \UV$，内阻 $r$ 约为 $0.5 \UO$），电压表（量程 $0 \sim 3 \UV$，内阻约为 $3 \UkO$），电流表（量程 $0 \sim 0.5 \UA$，内阻约为 $0.5 \UO$），滑动变阻器 $R$（最大阻值为 $20 \UO$），待测金属丝 $R_{x}$，刻度尺，螺旋测微器，开关 S，导线若干，实验电路如图~\ref{2026海南15a} 所示。
\threepicture[labelA=2026海南15a, labelB=2026海南15b, labelC=2026海南15c, h=2.8cm]
{figs/15a.png}{figs/15b.png}{figs/15c.png}

回答下列问题：
\begin{enumerate}
	\item
	按实验电路图（a）将图（b）的实物电路连接完整。
	\item
	用刻度尺测得接入电路的金属丝的有效长度 $L = 62.50 \Ucm$，用螺旋测微器测量金属丝不同位置的直径，某次测量的示数如图~\ref{2026海南15c} 所示，该读数 $d =$ \tkanswer{$0.500$} $\Umm$，三次测量后得到的平均值恰好与 $d$ 相等。
	\item
	闭合开关 S，改变滑动变阻器 R 滑片的位置，测得多组电压 $U$ 和电流 $I$ 的值如表 1 所示，根据表 1 的数据在图~\ref{2026海南15d} 的坐标纸上补齐数据点，并作出 $U-I$ 图线。

	\begin{center}
	\begin{tabular}{|c|c|c|c|c|c|c|}
	\hline
	$U/\UV$ & 0.99 & 1.26 & 1.51 & 1.74 & 1.99 & 2.25 \\
	\hline
	$I/\UA$ & 0.200 & 0.250 & 0.300 & 0.350 & 0.400 & 0.450 \\
	\hline
	\end{tabular}
	\end{center}
	\item
	由 $U-I$ 图像可求出所测金属丝的电阻 $R_{x} =$ \tkanswer{$5.0$} $\UO$，电阻率 $\rho =$ \tkanswer{$1.6\times10^{-6}$} $\UOm$（结果均保留 2 位有效数字）。
	\item
	如果某同学在连接实物电路时误接成如图~\ref{2026海南15e} 所示的电路。结合（4）中求得的金属丝 $R_{x}$ 阻值，若取电源内阻为 $0.5 \UO$，电流表内阻为 $0.5 \UO$，当滑动变阻器 $R$ 滑片滑到某一位置时，金属丝发热功率最小，该最小功率为 \tkanswer{$0.37$} $\UW$（结果保留 2 位有效数字）。
\end{enumerate}
\twopicture[labelA=2026海南15d, labelB=2026海南15e, numA=d, numB=e, h=4.0cm]
{figs/15d.png}{figs/15e.png}

%% answer:
\jdanswer{
\begin{enumerate}
	\item
	（如图）

	\item
	$0.500$

	\item
	（如图）

	\item
	$5.0$，$1.6\times10^{-6}$

	\item
	$0.37$
\end{enumerate}
}
%% memo:
\memoanswer{
（1）根据电路图可得实物连接图如图所示。

\onepicture[w=5.0cm]{figs/15_an1.png}

（2）螺旋测微器的读数为 $d = 0.5 \Umm + 0.0\times 0.01 \Umm = 0.500 \Umm$。

（3）根据描点法可得 $U-I$ 图线如图所示。

\onepicture[w=6.0cm]{figs/15_an2.png}

（4）根据 $U-I$ 图像可求出所测金属丝的电阻 $R_{x}=\frac{\Delta U}{\Delta I}\approx5.0 \UO$。

根据电阻定律 $R_{x}=\rho\frac{L}{S}$，其中 $S=\pi\left(\frac{d}{2}\right)^{2}$，代入数据可得 $\rho \approx 1.6 \times 10^{-6} \UOm$。

（5）根据题图可知图（e）电路中，滑动变阻器的滑片将自身分为两部分，并联后再与 $R_{x}$、电流表串联，设滑动变阻器一部分阻值为 $R_{p}$，另一部分阻值为 $R - R_{p}$，并联后等效阻值为 $\frac{(R - R_{p})R_{p}}{(R - R_{p}) + R_{p}}$，由数学知识可知当 $R - R_{p} = R_{p}$ 时，等效阻值最大，此时等效阻值为 $5 \UO$，电路中总电阻最大，流过 $R_{x}$ 的电流最小为 $I_{\min} = \frac{3}{5 + 5 + 0.5 + 0.5} \UA$。

金属丝发热功率最小为 $P_{\min} = I_{\min}^{2} R_{x}$，代入数据可得 $P_{\min} \approx 0.37 \UW$。
}

\item
%% number: 16
%% paperName: 2026年海南高考第 16 题 10 分
%% typeId: 6
%% type: 计算题
%% chapter: 运动和力
%% point: 动量定理
%% method: 无
%% score: 10
%% degree: 650
%% duplicateId: 0
%% body:
如图，一质量为 $3 \Ukg$ 的椰子从树上距水平地面 $5 \Um$ 高处由静止开始下落，与地面碰撞后静止。不计空气阻力，取重力加速度 $g = 10 \Umsq$。
\begin{enumerate}
	\item
	求椰子下落至地面所需要的时间；
	\item
	求椰子与地面碰撞前瞬间速度的大小；
	\item
	若碰撞时间为 $0.01 \Us$，碰撞过程中不计重力，求碰撞过程地面对椰子平均作用力的大小。
\end{enumerate}
\onepicture[w=4.5cm, align=right]{figs/16.png}

%% answer:
\jdanswer{
\begin{enumerate}
	\item
	$1 \Us$

	\item
	$10 \Ums$

	\item
	$3000 \UN$
\end{enumerate}
}
%% memo:
\memoanswer{
（1）椰子由静止开始自由下落，位移满足自由落体位移公式 $h = \frac{1}{2}gt^{2}$，整理得 $t=\sqrt{\frac{2h}{g}}=\sqrt{\frac{2\times5}{10}} \Us=1 \Us$。

（2）根据自由落体速度公式 $v = gt$，代入 $t = 1 \Us$，得 $v = 10 \Ums$。

（3）取竖直向上为正方向，根据动量定理 $F\Delta t = 0 - (-mv)$，整理得 $F=3000 \UN$。
}

\item
%% number: 17
%% paperName: 2026年海南高考第 17 题 12 分
%% typeId: 6
%% type: 计算题
%% chapter: 磁场
%% point: 带电粒子在磁场中的运动
%% method: 无
%% score: 12
%% degree: 450
%% duplicateId: 0
%% body:
如图，平行金属板 M、N 水平放置，$O_{1}O_{2}$ 是两板间的中心线，板间的矩形区域 Ⅰ 中存在竖直向下的匀强电场和垂直纸面向里的匀强磁场。金属板右侧有一圆心为 $O$、半径为 $R$ 的圆形区域 Ⅱ，$O$ 在 $O_{1}O_{2}$ 的延长线上，区域 Ⅱ 内有垂直纸面向外的匀强磁场。金属板左侧有粒子发射器 S 位于 $O_{1}O_{2}$ 上方 $\frac{\sqrt{3}}{2}R$ 处。发射器以初速度 $v_{0}$ 水平向右射出带正电的粒子，粒子在区域 Ⅰ 恰好做匀速直线运动，从 $A$ 点射入区域 Ⅱ 后从 $C$ 点射出。已知 $A$、$C$ 两点关于 $O_{1}O_{2}$ 对称，金属板间电场强度为 $E$，粒子质量为 $m$、电荷量为 $q$，不计粒子重力。
\begin{enumerate}
	\item
	求区域 Ⅰ 中磁场磁感应强度的大小；
	\item
	求区域 Ⅱ 中磁场磁感应强度的大小和粒子在区域 Ⅱ 中的运动时间；
	\item
	调整区域 Ⅱ 中磁场磁感应强度的大小后，从 $A$ 点入射的粒子能从区域 Ⅱ 与圆心 $O$ 等高的 $D$ 点射出，求调整后区域 Ⅱ 中磁感应强度的大小。
\end{enumerate}
\onepicture[w=9.0cm, align=right]{figs/17.png}

%% answer:
\jdanswer{
\begin{enumerate}
	\item
	$B_{1} = \frac{E}{v_{0}}$

	\item
	$B_{2} = \frac{2\sqrt{3}mv_{0}}{3qR}$；$t = \frac{\sqrt{3}\pi R}{2v_{0}}$

	\item
	$B_{2}' = \frac{\sqrt{3}mv_{0}}{3qR}$
\end{enumerate}
}
%% memo:
\memoanswer{
（1）粒子在区域 Ⅰ 做匀速直线运动，竖直方向受力平衡：电场力与洛伦兹力大小相等。带正电粒子受竖直向下的电场力 $qE$，由左手定则，垂直纸面向里的磁场对水平向右运动的粒子的洛伦兹力竖直向上，平衡条件为 $qE = qv_{0}B_{1}$，整理得区域 Ⅰ 磁感应强度 $B_{1}=\frac{E}{v_{0}}$。

（2）入射点 A 到中心线 $O_{1}O_{2}$ 的竖直距离为 $\frac{\sqrt{3}}{2}R$，圆形区域 Ⅱ 的圆心 O 在中心线上，因此 $OA$ 与水平中心线的夹角满足 $\sin\theta=\frac{\frac{\sqrt{3}R}{2}}{R}=\frac{\sqrt{3}}{2}$，即 $\theta=60^{\circ}$。

由于 A、C 关于 $O_{1}O_{2}$ 对称，粒子在区域 Ⅱ 中轨迹对应的圆心角为 $180^{\circ}$，由几何关系可得粒子做圆周运动的轨道半径 $r_{1}=R\sin 60^{\circ}=\frac{\sqrt{3}R}{2}$。

洛伦兹力提供向心力：$qv_{0}B_{2}=m\frac{v_{0}^{2}}{r_{1}}$，得 $B_{2}=\frac{2\sqrt{3}mv_{0}}{3qR}$。

粒子在磁场中运动的周期 $T=\frac{2\pi m}{qB_{2}}$，轨迹圆心角为 $\pi$，因此运动时间 $t=\frac{1}{2}T=\frac{\sqrt{3}\pi R}{2v_{0}}$。

（3）调整后粒子从 A 点水平向右入射，从与 O 等高的 D 点射出区域 Ⅱ，由几何关系 $\left(r_{2}-\frac{\sqrt{3}}{2}R\right)^{2}+\left(R+\frac{\sqrt{3}}{2}R\tan 30^{\circ}\right)^{2}=r_{2}^{2}$，得轨道半径 $r_{2}=\sqrt{3}R$。

\onepicture[w=6.5cm]{figs/17_an.png}

代入洛伦兹向心力公式 $qv_{0}B_{2}'=m\frac{v_{0}^{2}}{r_{2}}$，得 $B_{2}'=\frac{\sqrt{3}mv_{0}}{3qR}$。
}

\item
%% number: 18
%% paperName: 2026年海南高考第 18 题 16 分
%% typeId: 6
%% type: 计算题
%% chapter: 运动和力
%% point: 传送带
%% method: 无
%% score: 16
%% degree: 350
%% duplicateId: 0
%% body:
如图，轻质弹性绳左端固定于 $O$ 点，右端与一置于水平地面上的木板 A 相连。连接点与 $O$ 点等高。A 上左端有一带电物体 B，空间存在水平向右、场强大小为 $5\times10^{4} \UVm$ 的匀强电场 $E$，已知地面光滑，A、B 之间的动摩擦因数为 0.75。弹性绳原长为 $0.5 \Um$。弹力大小与伸长量满足胡克定律。劲度系数为 $12.5 \UNm$。A 和 B 的质量均为 $1 \Ukg$，A 不带电，B 的电荷量为 $+2\times10^{-4} \UC$。$t = 0$ 时，A 的左端到 $O$ 的水平距离为 $0.4 \Um$。A 和 B 由静止开始释放。设木板 A 足够长，A、B 之间的最大静摩擦力等于滑动摩擦力。运动过程中 B 的电荷量保持不变，取重力加速度 $g = 10 \Umsq$。
\begin{enumerate}
	\item
	求 $t = 0$ 时 A 和 B 的共同加速度大小；
	\item
	求 B 与 A 开始发生相对滑动时 B 的速度大小；
	\item
	当 B 的速度为 $3 \Ums$ 时剪断绳右端，又经 $0.4 \Us$ 后 A 和 B 速度相等，求从 $t = 0$ 开始 A、B 之间因摩擦产生的热量。
\end{enumerate}
\onepicture[w=6.0cm, align=right]{figs/18.png}

%% answer:
\jdanswer{
\begin{enumerate}
	\item
	$5 \Umsq$

	\item
	$2 \Ums$

	\item
	$9 - \frac{3\sqrt{7}}{2} \UJ$
\end{enumerate}
}
%% memo:
\memoanswer{
（1）B 受到的电场力：$F_{\text{电}}=qE=10 \UN$。

对 A、B 整体由牛顿第二定律：$F_{\text{电}}=2ma$，解得 $a = 5 \Umsq$。

（2）A、B 刚好发生相对滑动时，对 B 有 $qE - \mu mg = ma_{0}$，得此时共同临界加速度 $a_{0}=2.5 \Umsq$；设此时弹性绳的伸长量为 $x$，对 A 有 $\mu mg - kx = ma_{0}$，解得 $x = 0.4 \Um$。

对应 B 向右的位移 $s = L_{0} - 0.4 \Um + x = 0.5 \Um$，此过程中由动能定理 $F_{\text{电}}s - \frac{1}{2}kx^{2} = \frac{1}{2}(m + m)v^{2}$，代入数据解得 $v = 2 \Ums$。

（3）发生相对滑动后，B 的加速度 $a_{\mathrm{B}}=\frac{F_{\text{电}}-\mu m_{\mathrm{B}}g}{m_{\mathrm{B}}}=2.5 \Umsq$。

剪断绳后，A 的加速度 $a_{A} = \mu g = 7.5 \Umsq$；经 $t_{0}=0.4 \Us$ 共速，共速速度 $v_{\text{共}}=v_{\mathrm{B}}+a_{\mathrm{B}}t_{0}=4 \Ums$，则剪断绳时 A 的速度 $v_{A}=v_{\text{共}}-a_{A}t_{0}=1 \Ums$。

由刚好发生相对滑动到 B 的速度达到 $v_{\mathrm{B}}=3 \Ums$ 过程：此过程 B 的位移 $x_{\mathrm{B1}}=\frac{v_{\mathrm{B}}^{2}-v^{2}}{2a_{\mathrm{B}}}=1 \Um$。

对 A，根据动能定理 $\mu mgx_{\mathrm{A1}}-\left[\frac{1}{2}k\left(x_{\mathrm{A1}}+x\right)^{2}-\frac{1}{2}kx^{2}\right]=\frac{1}{2}mv_{\mathrm{A}}^{2}-\frac{1}{2}mv^{2}$，解得 A 的位移 $x_{\mathrm{A1}}=\frac{1+\sqrt{7}}{5} \Um$。

相对位移 $\Delta x_{1}=x_{\mathrm{B1}}-x_{\mathrm{A1}}=\frac{4-\sqrt{7}}{5} \Um$；由剪断绳到共速过程相对位移 $\Delta x_{2}=x_{\mathrm{B2}}-x_{\mathrm{A2}}=\frac{v_{\mathrm{B}}+v_{\text{共}}}{2}t_{0}-\frac{v_{A}+v_{\text{共}}}{2}t_{0}=0.4 \Um$。

从 $t = 0$ 开始 A、B 之间因摩擦产生的热量 $Q = \mu mg(\Delta x_{1} + \Delta x_{2}) = 9 - \frac{3\sqrt{7}}{2} \UJ$。
}

\end{enumerate}

\end{document}
